% !TeX root = thesis.tex

\chapter{Related work}
In the previous chapter I have stressed the paramount importance of periodically integrating one's changes into the upstream repository. Additionally, Continuous Integration was introduced as both a practice and a tool to facilitate this often complex and time-consuming process. However, Continuous Integration is not the golden bullet for software engineering. In this chapter I will investigate the flip side of applying this practice. After every integration, all of the unit and regression tests in the test suite must be executed to ensure that the integration was successful and that no new bugs have been introduced. As the project evolves, the size of the codebase increases and consequently the amount of tests will increase as well in order to maintain a sufficiently high coverage level. An increase in the size of the test suite will inevitably lead to an increase in test duration \cite{evaluationoftestsuiteminimization}, which imposes an issue of scaling. Walcott, Soffa and Kapfhammer illustrate the magnitude of this problem by providing an example of a codebase consisting of 20.000 lines, for which the tests require up to seven weeks to complete \cite{10.1145/1146238.1146240}.\\

\noindent Fortunately, multiple developers and researchers have found some techniques that can be used to address the scalability issues of growing test suites. The techniques currently known to literature can be classified in three categories. Developers can either apply \emph{\tsm{}}, \emph{\tcs{}} or \emph{\tcp{}} \cite{evaluationoftestsuiteminimization}. All three techniques are applicable to any test suite, however there is a trade-off to be made. Depending on which technique is chosen, it will either have a major impact on the duration of the test suite execution in exchange for a reduced test coverage level, or it will result in a higher test adequacy.\\

\noindent In the following sections I will first elaborate on the details of these three approaches, then I will provide accompanying algorithms that can be used. Since the approaches share common ideas, the algorithms can (albeit with minor modifications) be reused across all approaches. In the final section I will illustrate some examples of the discussed techniques and algorithms that are currently integrated in existing software testing frameworks.

% !TeX root = ../thesis.tex

\section{Classification of approaches}

% !TeX root = ../../thesis.tex

\subsection{\tsm{}}
\label{ssec:tsm}
The first technique is called \acrfull{tsm}, also referred to as \emph{Test Suite Reduction} in literature. This technique will try to reduce the size of the test suite by permanently removing redundant test cases. This problem has been formally defined by Rothermel \cite{10.1002/stv.430} in \cref{def:tsm} and illustrated in \Cref{fig:tsm}.

\begin{definition}[\tsm{}]
\label{def:tsm}
\mbox{}\\Given:
\begin{itemize}
	\item $T = \{t_1, \dots, t_n\}$ a test suite consisting of test cases $t_j$.
	\item $R = \{r_1, \dots, r_m\}$ a set of requirements that must be satisfied in order to provide the desired ``adequate'' testing of the program.
	\item $\{T_1, \dots, T_m\}$ subsets of test cases in $T$, one associated with each of the requirements $r_i$, such that any one of the test cases $t_j \in T_i$ can be used to satisfy requirement $r_i$.
\end{itemize}

\noindent Subsequently, we can define \tsm{} as the task of finding a subset $T'$ of test cases $t_j \in T$ that satisfies every requirement $r_i$.
\end{definition}

\noindent If we apply the concepts of the previous chapter to the above definition, we can interpret the set of requirements $R$ as source code lines that must be covered. A requirement $r_i$ can subsequently be satisfied by any test case $t_j \in T$ that belongs to the subset $T_i$. Observe that the problem of finding $T'$ is closely related to the \emph{hitting set problem} (\cref{def:hitting-set}) \cite{10.1002/stv.430}.

\begin{definition}[Hitting Set Problem]
\label{def:hitting-set}
\mbox{}\\Given:
\begin{itemize}
	\item $S = \{s_1, \dots, s_n\}$ a finite set of elements.
	\item $C = \{c_1, \dots, c_n\}$ a collection of sets, with $\forall c_i \in C : c_i \subseteq S$.
	\item $K$ a positive integer, $K \le |S|$.
\end{itemize}

\noindent The hitting set is a subset $S' \subseteq S$ such that $S'$ contains at least one element from each subset in $C$.
\end{definition}

\noindent In the context of \tsm{}, $T'$ corresponds to the hitting set of $T_i$s. In order to effectively minimise the amount of tests in the test suite, $T'$ should be the minimal hitting set \cite{10.1002/stv.430}. Since we can reduce this problem to the NP-complete \emph{Vertex Cover}-problem, we know that this problem is NP-complete as well \cite{10.5555/574848}.

\begin{figure}[htbp!]
	\centering
	\includegraphics[width=\textwidth]{assets/tikz/approach-tsm.tikz}
	\caption{\tsm{}}
	\label{fig:tsm}
\end{figure}
% !TeX root = ../../thesis.tex

\subsection{\tcs{}}
The second algorithm closely resembles the previous one. Instead of determining the minimal hitting set of the test suite in order to permanently remove tests, this algorithm has a notion of context. Prior to the execution of the tests, the algorithm performs a \emph{white-box static analysis} of the codebase to identify which parts have been changed. Subsequently, only the tests regarding modified parts are executed, making the selection temporary (\autoref{fig:tcs}) and modification-aware \cite{10.1002/stv.430}. Rothermel and Harrold define this formally in \autoref{def:tcs}.

\begin{definition}[\tcs{}]
\label{def:tcs}
\mbox{}\\Given:
\begin{itemize}
	\item $P$ the previous version of the codebase
	\item $P'$ the current (modified) version of the codebase
	\item $T$ the test suite
\end{itemize}

\noindent \tcs{} aims to find a subset $T' \subseteq T$ that is used to test $P'$. 
\end{definition}

\begin{figure}[htbp!]
	\centering
	\includegraphics[width=0.45\textwidth]{assets/placeholder.pdf}
	\caption{\tcs{}}
	\label{fig:tcs}
\end{figure}
% !TeX root = ../../thesis.tex

\subsection{\tcp{}}
Both \acrshort{tsm} and \acrshort{tcs} attempt to execute as few tests as possible to reduce the execution time of the test suite. Nevertheless, in some cases, we may require to execute every test case to guarantee correctness. In this situation, we can still optimise the test suite. \acrfull{tcp} aims to find a permutation of the sequence of test cases, rather than eliminating specific tests from being executed (\Cref{fig:tcp}). We choose the order of the permutation in such a way that we can complete a predefined objective as soon as possible. Once we have achieved our objective, we can early terminate the execution of the test suite. In the worst-case scenario, we will still execute every test case. Some examples of objectives include covering as many lines of code as fast as possible or executing tests ordered on their probability of failure \cite{10.1002/stv.430}. \Cref{def:tcp} provides a formal definition of this approach.

\begin{definition}[\tcp{}]
\label{def:tcp}
\mbox{}\\Given:
\begin{itemize}
	\item $T$ the test suite
	\item $PT$ the set of permutations of $T$
	\item $f: PT \mapsto \mathbb{R}$ a function from a subset to a real number, this function is used to compare sequences of test cases to find the optimal permutation.
\end{itemize}

\noindent \tcp{} finds a permutation $T' \in PT$ such that $\forall T'' \in PT : f(T') \ge f(T'') \Rightarrow (T'' \ne T')$ 
\end{definition}

\begin{figure}[htbp!]
	\centering
	\includegraphics[width=\textwidth]{assets/tikz/approach-tcp.tikz}
	\caption{\tcp{}}
	\label{fig:tcp}
\end{figure}
\section{Algorithms}
\noindent This thesis prefers to use TCP since this technique does not incur the risk of false-negative failing test cases. In order to determine the optimal order of execution, this thesis presents three existing algorithms.\\

\noindent The input data for these algorithms is threefold:\\

\begin{itemize}
\item \textbf{Affected test cases:} By combining previous coverage results and the list of changes that the developer has made to the code, the framework can estimate which test cases are likely affected by those changes and assign a higher priority.

\item \textbf{Historical data:} Next, historical failure data can be used. Suppose that a test case has failed in its previous execution, then there exists a chance that it will subsequently fail in the current run.

\item \textbf{Execution timings:} Finally, if two test cases are considered equally likely to fail, the average duration of the test case can be used as a tie-breaker. Since the objective of TCP is to optimise the test suite, the test case with the lowest duration should be preferred.
\end{itemize}

\subsection{Greedy algorithm}
\noindent The first algorithm is a greedy heuristic that was initially designed as an algorithm for the set-covering problem \cite{evaluationoftestsuiteminimization}. This heuristic starts with an empty set of test cases and the set of all code lines $C$. Next, the algorithm iteratively selects the test case that contributes the most code lines that are not yet covered, updating $C$ after every selected test case. The algorithm halts when either all test cases are selected or $C$ is empty. In order to modify this algorithm to make it applicable to TCP, the selection order must be preserved and used as the prioritised sequence.

\subsection{HGS algorithm}
\noindent The second algorithm was created by Harrold, Gupta and Soffa \cite{hgs}. As opposed to the greedy heuristic, this algorithm uses a different perspective. First, the algorithm sorts the code lines increasingly based on the number of test cases that cover them. The motivation for this sorting operation is that some test cases must inevitably be executed, as they are the only test cases that cover a given set of lines. However, these test cases can also cover other lines and therefore make other test cases redundant. The algorithm iterates over the code lines in this order and selects one corresponding test case in each iteration. Afterwards, the order is updated to remove source code lines which are now covered by the selected test case. This process is repeated until there are no lines left to cover.

\clearpage

\subsection{ROCKET algorithm}
\noindent Finally, this thesis considers the ROCKET algorithm \cite{6676952}. This algorithm prioritises test cases by assigning a score to every test case, which is calculated using historical failure data. Afterwards, the algorithm computes the cumulative score $CS_t$ of every test case $t$ and defines the following objective function $g$, in which $E_t$ represents the execution time of the test case:
$$g = (maximise(CS_t), minimise(E_t))$$

\noindent Finally, the algorithm optimises this function to determine the ideal order of execution $S$, as follows:
$$(\forall i \in 1 \dots n)(g(S_i) \ge g(S_{i+1})$$

\section{Existing implementations}
- OpenClover (enkel Java) heeft hier misschien support voor